% 2.2.3
% Ficheiro incluído por main.tex

\subsection[Configuração do Samba (smb.conf)]{Configuração do Samba (\texttt{smb.conf})}

\esquema{2-2-3-smb-conf}{O smb.conf define o Ubuntu como membro do domínio ISEC}

\passo{Passo 1 - Guardar a configuração original}

\begin{lstlisting}[style=bashstyle]
sudo cp -a /etc/samba/smb.conf /etc/samba/smb.conf.bak
\end{lstlisting}

\passo{Passo 2 - Definir a secção \texttt{[global]}}

\begin{lstlisting}[style=bashstyle]
sudo nano /etc/samba/smb.conf
\end{lstlisting}

Substituir toda a secção \texttt{[global]} (até à linha \texttt{\#======================= Share Definitions =======================}) por:

\begin{lstlisting}[style=cfgstyle]
[global]
    # Identificação do domínio Active Directory
    workgroup = ISEC
    realm = ISEC.LOCAL
    security = ADS

    # Nome NetBIOS (máximo 15 caracteres) - nome da máquina no AD
    netbios name = UBU-WINBIND

    # Mapeamento de identificadores Windows para Linux
    idmap config * : backend = tdb
    idmap config * : range = 3000-7999
    idmap config ISEC : backend = rid
    idmap config ISEC : range = 10000-999999

    # Integração com Winbind
    winbind use default domain = yes
    winbind enum users = yes
    winbind enum groups = yes
    winbind refresh tickets = yes

    # Ambiente dos utilizadores do domínio
    template shell = /bin/bash
    template homedir = /home/%U

    # Integração com PAM
    obey pam restrictions = yes
\end{lstlisting}

Guardar com \texttt{Ctrl+O}, \texttt{Enter}; sair com \texttt{Ctrl+X}.

\begin{center}\small
\begin{tabularx}{\linewidth}{|>{\raggedright\arraybackslash}p{0.32\linewidth}|>{\raggedright\arraybackslash}X|}
\hline
\rowcolor{azul!10}\textbf{Parâmetro} & \textbf{O que faz} \\
\hline
\texttt{workgroup}, \texttt{realm}, \texttt{security = ADS} & Domínio NetBIOS \texttt{ISEC}, realm Kerberos \texttt{ISEC.LOCAL} e modo membro de Active Directory \\
\hline
\texttt{netbios name = UBU-WINBIND} & Nome da máquina no AD. É diferente do cliente SSSD (\texttt{UBUCLIENT}), por isso os dois podem estar no mesmo domínio \\
\hline
\texttt{idmap config * : ...} (\texttt{tdb}) & UIDs/GIDs para contas fora do domínio (ex.: \texttt{BUILTIN}), entre 3000 e 7999 \\
\hline
\texttt{idmap config ISEC : ...} (\texttt{rid}) & UIDs/GIDs dos utilizadores do domínio calculados a partir do RID do AD (10000 + RID): são sempre os mesmos \\
\hline
\texttt{winbind use default domain = yes} & Os utilizadores escrevem-se sem domínio: \texttt{professor.teste} em vez de \texttt{ISEC\textbackslash{}professor.teste} \\
\hline
\texttt{winbind enum users/\allowbreak{}groups = yes} & Permite listar todos os utilizadores e grupos (\texttt{wbinfo -u}, \texttt{wbinfo -g}) \\
\hline
\texttt{winbind refresh tickets = yes} & Renova automaticamente os tickets Kerberos dos utilizadores \\
\hline
\texttt{template shell}, \texttt{template homedir} & Shell e pasta pessoal dos utilizadores do domínio (\texttt{/\allowbreak{}home/\allowbreak{}professor.teste}) \\
\hline
\texttt{obey pam restrictions = yes} & O Samba respeita as regras do PAM \\
\hline
\end{tabularx}
\end{center}

\passo{Passo 3 - Verificar e validar a sintaxe}

\begin{lstlisting}[style=bashstyle]
sudo testparm -s
\end{lstlisting}

\emph{Verificação esperada:} \texttt{Loaded services file OK.} e \textbf{\texttt{Server role: ROLE\_DOMAIN\_MEMBER}}, seguido dos parâmetros da \texttt{[global]}. O aviso \texttt{Weak crypto is allowed by GnuTLS} é normal.

\problemas{\hyperref[ec-b2]{Erros Comuns - 2.2 Winbind}}
